\newpage
\chapter{PENDAHULUAN} \label{Bab I}

\section{Latar Belakang Masalah} \label{I.Latar Belakang}
Kemajuan teknologi informasi dan komunikasi telah mengubah cara masyarakat memproduksi, mendistribusikan, dan mengonsumsi informasi digital. Dalam konteks tersebut, media visual seperti citra dan video memiliki peran penting karena sering dipersepsikan sebagai representasi yang dekat dengan realitas. Namun, perkembangan kecerdasan buatan generatif juga meningkatkan risiko manipulasi informasi visual dalam skala besar. \textit{Global Risks Report 2026} yang diterbitkan oleh \textit{World Economic Forum} menunjukkan bahwa misinformasi dan disinformasi merupakan ancaman yang konsisten lintas periode proyeksi, yaitu menempati peringkat kelima pada tahun 2026, peringkat kedua dalam proyeksi dua tahun hingga 2028, dan peringkat keempat dalam proyeksi sepuluh tahun hingga 2036 \cite{WorldEconomicForumWEF2026}. Salah satu bentuk manipulasi visual yang menonjol dalam konteks ini adalah \textit{deepfake}, yaitu media sintetis berbasis \textit{deep learning} yang mampu menghasilkan citra atau video wajah yang tampak realistis, sehingga berpotensi digunakan untuk penipuan, manipulasi opini, penyalahgunaan identitas digital, dan penyebaran informasi palsu \cite{Tolosana2020,Mirsky_2021,BhaskardeepKhaund2025}. \par

Ancaman \textit{deepfake} tidak lagi terbatas pada ranah hiburan atau eksperimen teknis, tetapi telah berkembang menjadi persoalan keamanan informasi dan kepercayaan digital. Tolosana dkk. menjelaskan bahwa kemajuan teknik manipulasi wajah, khususnya yang ditopang oleh \textit{deep learning} dan \textit{Generative Adversarial Network}, telah menghasilkan konten palsu yang semakin meyakinkan serta menimbulkan tantangan serius bagi sistem deteksi \cite{Tolosana2020}. Sejalan dengan itu, Khaund menunjukkan bahwa media sintetis modern dapat melemahkan mekanisme autentikasi yang bergantung pada biometrik dan verifikasi dokumen, sehingga \textit{deepfake} menjadi ancaman yang relevan bagi ekosistem identitas digital \cite{BhaskardeepKhaund2025}. Kondisi tersebut menegaskan bahwa pengembangan sistem deteksi otomatis tetap diperlukan, terutama untuk menghadapi konten video yang beredar pada lingkungan dunia nyata dengan kualitas visual yang beragam. \par

Dalam perkembangannya, pendekatan deteksi \textit{deepfake} dapat dibedakan menjadi pendekatan spasial dan pendekatan spasiotemporal. Pendekatan spasial bekerja pada tingkat \textit{frame} tunggal dan berfokus pada artefak visual statis, seperti ketidakwajaran tekstur, batas wajah, atau inkonsistensi detail lokal \cite{Tolosana2020,Cao2022}. Pendekatan ini relatif sederhana karena setiap \textit{frame} dapat dianalisis sebagai citra individual. Selain itu, pendekatan berbasis Vision Transformer juga menunjukkan potensi dalam mempelajari hubungan global antarbagian citra, sehingga relevan untuk mendeteksi artefak visual yang tersebar pada area wajah \cite{dosovitskiy2021imageworth16x16words,app12062953}. Meskipun demikian, pendekatan spasial memiliki keterbatasan karena tidak memodelkan hubungan antar-\textit{frame} secara eksplisit. Akibatnya, jika jejak manipulasi hanya muncul melalui perubahan gerakan atau ketidakkonsistenan temporal, pendekatan spasial berpotensi tidak menangkap informasi tersebut secara optimal. \par

Sebagai alternatif, pendekatan spasiotemporal memanfaatkan informasi pada tingkat video dengan memperhatikan hubungan antar-\textit{frame}. Gu dkk. menunjukkan bahwa deteksi video \textit{deepfake} tidak hanya dapat dilakukan melalui analisis \textit{frame} tunggal, tetapi juga melalui pemodelan inkonsistensi spasial dan temporal pada video hasil manipulasi \cite{Gu2021}. Haliassos dkk. juga menunjukkan bahwa representasi gerakan mulut yang bersifat spasiotemporal dapat membantu meningkatkan generalisasi deteksi \textit{face forgery} pada berbagai kondisi distorsi \cite{Haliassos2021}. Meskipun demikian, pendekatan spasiotemporal tetap perlu diuji secara empiris karena efektivitas informasi temporal dapat dipengaruhi oleh karakteristik input video, seperti jumlah \textit{frame} yang digunakan, jarak antar-\textit{frame} yang dipilih, serta kualitas spasial \textit{frame}. Oleh karena itu, penelitian ini membandingkan pendekatan spasial dan spasiotemporal melalui konfigurasi input yang sesuai dengan eksperimen yang dilakukan, yaitu variasi resolusi, jumlah \textit{frame}, dan \textit{stride}. \par

Permasalahan tersebut menjadi semakin penting pada skenario dunia nyata (\textit{in-the-wild}). Zi dkk. memperkenalkan \textit{WildDeepfake} sebagai \textit{dataset} \textit{deepfake} dunia nyata yang dikumpulkan dari internet dan terdiri atas 7.314 video wajah \cite{Zi2024}. \textit{Dataset} ini memiliki variasi kualitas visual, pose, pencahayaan, dan panjang video yang beragam. Zi dkk. juga menunjukkan bahwa \textit{WildDeepfake} lebih menantang dibandingkan \textit{dataset} yang disusun dalam kondisi terkontrol, karena performa beberapa detektor dasar mengalami penurunan ketika diuji pada \textit{dataset} ini \cite{Zi2024}. Dengan demikian, \textit{WildDeepfake} menjadi konteks yang relevan untuk menguji apakah pendekatan spasiotemporal benar-benar memberikan keuntungan dibandingkan pendekatan spasial, terutama ketika konfigurasi input seperti resolusi, jumlah \textit{frame}, dan \textit{stride} diubah secara terkontrol. \par

Berdasarkan kondisi tersebut, penelitian ini berfokus pada analisis komparatif antara pendekatan spasial dan spasiotemporal untuk klasifikasi video \textit{deepfake}. Pendekatan spasial direpresentasikan oleh Vision Transformer (ViT), sedangkan pendekatan spasiotemporal direpresentasikan oleh VideoMAE. Kedua model tersebut dipilih karena sama-sama berasal dari keluarga Transformer, tetapi memiliki cara kerja yang berbeda dalam memproses data video. ViT memproses \textit{frame} sebagai citra tunggal melalui urutan \textit{image patches}, sedangkan VideoMAE memproses klip video dengan mempertimbangkan hubungan spasial dan temporal secara bersamaan \cite{dosovitskiy2021imageworth16x16words,tong2022videomaemaskedautoencodersdataefficient}. Dengan rancangan tersebut, penelitian ini tidak berangkat dari asumsi bahwa pendekatan spasiotemporal pasti lebih unggul, melainkan menguji secara empiris sejauh mana informasi temporal benar-benar memberikan kontribusi terhadap klasifikasi video \textit{deepfake} pada data \textit{in-the-wild}. \par

Penelitian ini menggunakan \textit{WildDeepfake} sebagai \textit{dataset} utama untuk pelatihan, validasi, dan pengujian internal. Selain itu, evaluasi lintas-\textit{dataset} secara \textit{zero-shot} pada \textit{Celeb-DF v2} digunakan sebagai pengujian tambahan untuk melihat kemampuan generalisasi model pada distribusi data eksternal yang berbeda dari \textit{dataset} pelatihan \cite{li2020celebdflargescalechallengingdataset}. Untuk itu, penelitian ini tidak hanya membandingkan efektivitas pendekatan spasial dan spasiotemporal, tetapi juga menganalisis pengaruh konfigurasi pelatihan serta konfigurasi input spasial dan temporal, seperti resolusi, jumlah \textit{frame}, dan \textit{stride}, terhadap performa klasifikasi video \textit{deepfake}. Dengan demikian, penelitian ini diharapkan dapat memberikan gambaran yang lebih terarah mengenai pendekatan yang lebih efektif dan lebih tangguh ketika diterapkan pada data dunia nyata maupun pada skenario generalisasi lintas-\textit{dataset}. \par

\section{Rumusan Masalah} \label{I.Rumusan Masalah}
\begin{enumerate}[noitemsep]
	\item Bagaimana merancang dan mengimplementasikan model klasifikasi video \textit{deepfake} menggunakan pendekatan spasial berbasis Vision Transformer (ViT) dan pendekatan spasiotemporal berbasis VideoMAE pada dataset WildDeepfake?
	\item Bagaimana mengukur serta membandingkan performa pendekatan spasial berbasis Vision Transformer (ViT) dan pendekatan spasiotemporal berbasis VideoMAE berdasarkan \textit{accuracy}, \textit{precision}, \textit{recall}, \textit{F1-score}, AUC-ROC, dan \textit{confusion matrix}?
	\item Pendekatan manakah di antara model spasial berbasis Vision Transformer (ViT) dan model spasiotemporal berbasis VideoMAE yang paling efektif untuk klasifikasi video \textit{deepfake} pada dataset WildDeepfake?
\end{enumerate}

\section{Tujuan Penelitian} \label{I.Tujuan}
Berdasarkan rumusan masalah yang telah diuraikan, tujuan penelitian ini adalah sebagai berikut. \par

\begin{enumerate}[noitemsep]
	\item Merancang dan mengimplementasikan model klasifikasi video \textit{deepfake} menggunakan pendekatan spasial berbasis Vision Transformer (ViT) dan pendekatan spasiotemporal berbasis VideoMAE pada dataset WildDeepfake.
	\item Mengukur serta membandingkan performa pendekatan spasial berbasis Vision Transformer (ViT) dan pendekatan spasiotemporal berbasis VideoMAE berdasarkan \textit{accuracy}, \textit{precision}, \textit{recall}, \textit{F1-score}, AUC-ROC, dan \textit{confusion matrix}.
	\item Menentukan pendekatan yang paling efektif di antara model spasial berbasis Vision Transformer (ViT) dan model spasiotemporal berbasis VideoMAE untuk klasifikasi video \textit{deepfake} pada dataset WildDeepfake.
\end{enumerate}

\section{Batasan Masalah} \label{I.Batasan}
Batasan masalah pada penelitian ini adalah sebagai berikut. \par

\begin{enumerate}[noitemsep]
	\item \textit{WildDeepfake} digunakan sebagai \textit{dataset} utama untuk pelatihan, validasi, dan pengujian internal, sedangkan \textit{Celeb-DF v2} digunakan hanya sebagai \textit{dataset} evaluasi eksternal secara \textit{zero-shot}. Pada \textit{Celeb-DF v2} tidak dilakukan pelatihan ulang maupun \textit{fine-tuning}, dan pengujian lintas-\textit{dataset} dalam penelitian ini dibatasi hanya pada \textit{Celeb-DF v2}.
	\item Penelitian ini dibatasi pada klasifikasi biner untuk membedakan video ke dalam kelas \textit{real} dan \textit{fake}.
	\item Data yang digunakan berupa video wajah yang telah melalui proses \textit{preprocessing} awal pada dataset sumber, sehingga penelitian ini tidak mencakup tahap deteksi wajah, ekstraksi fitur wajah, maupun \textit{face alignment} dari video mentah.
	\item Model yang digunakan dibatasi pada arsitektur Vision Transformer (ViT) untuk pendekatan spasial dan VideoMAE untuk pendekatan spasiotemporal, dengan memanfaatkan (\textit{pretrained weights}) untuk proses \textit{fine-tuning}.
	\item Evaluasi kinerja model dibatasi pada metrik \textit{accuracy}, \textit{precision}, \textit{recall}, \textit{F1-score}, AUC-ROC, dan \textit{confusion matrix}.
	\item Penelitian ini tidak mencakup perancangan atau pengembangan aplikasi maupun antarmuka pengguna, serta tidak membahas implementasi sistem dalam skenario \textit{real-time} dan integrasi dengan data audio atau teks.
\end{enumerate}

\section{Manfaat Penelitian} \label{I.Manfaat}
Manfaat penelitian ini adalah sebagai berikut. \par

\begin{enumerate}[noitemsep]
	\item Menghasilkan implementasi dan evaluasi komparatif model klasifikasi video \textit{deepfake} menggunakan pendekatan spasial berbasis Vision Transformer (ViT) dan pendekatan spasiotemporal berbasis VideoMAE pada dataset WildDeepfake.
	\item Memberikan kontribusi ilmiah berupa analisis efektivitas relatif informasi spasial dan temporal dalam klasifikasi video \textit{deepfake} pada kondisi dunia nyata (\textit{in-the-wild}).
	\item Menjadi referensi bagi penelitian selanjutnya dalam pengembangan sistem klasifikasi atau deteksi \textit{deepfake} yang lebih \textit{robust}, khususnya pada data video dengan kualitas visual yang beragam dan terkompresi.
\end{enumerate}


\section{Sistematika Penulisan} \label{I.Sistematika}
Sistematika penulisan yang diterapkan dalam menyusun penelitian ini adalah sebagai berikut. \par

\subsection{Bab I Pendahuluan}
Pada Bab I, penulis menjelaskan latar belakang permasalahan, rumusan masalah, tujuan penelitian, batasan masalah, manfaat penelitian, dan sistematika penulisan. \par

\subsection{Bab II Tinjauan Pustaka}
Pada Bab II, penulis memberikan penjabaran tentang penelitian sebelumnya yang terkait dengan penelitian ini. Penulis juga memaparkan dasar teori yang berkaitan dengan penelitian yang dilakukan. \par

\subsection{Bab III Metode Penelitian}
Pada Bab III, penulis menjelaskan mengenai metodologi yang akan digunakan pada penelitian ini. Metodologi yang digunakan dapat dijelaskan dengan diagram alir yang dilanjutkan dengan penjelasan setiap tahapnya dan juga menjelaskan tentang perangkat apa yang digunakan oleh penulis untuk penelitian. \par

\subsection{Bab IV Hasil dan Pembahasan}
Pada Bab IV, penulis menjelaskan hasil dari penelitian yang didapatkan, dimulai dengan pengumpulan data, pengolahan data, hasil penelitian, dan pembahasan penelitian. \par

\subsection{Bab V Kesimpulan dan Saran}
Pada Bab V, penulis memberikan kesimpulan dari hasil penelitian yang didapatkan pada Bab IV, serta saran yang dapat dilakukan untuk penelitian selanjutnya.
